\documentclass[11pt]{amsart}
\ifdefined\XeTeXrevision\else\ifdefined\pdfoutput\pdfoutput=1\fi\fi

\usepackage[margin=1.2in]{geometry}
\usepackage{amsmath,amssymb,mathtools}
\usepackage{booktabs}
\usepackage{array}
\usepackage{xcolor}
\definecolor{linkblue}{RGB}{25,60,120}
\usepackage[colorlinks=true,linkcolor=linkblue,citecolor=linkblue,urlcolor=linkblue]{hyperref}
\hypersetup{
  pdftitle={Peaceable queens on the 15x15 and 16x16 tori: t(15)=20 and t(16)=32},
  pdfauthor={Jan Philipp Harries},
  pdfsubject={Peaceable queens on toroidal chessboards},
  pdfkeywords={peaceable queens, toroidal chessboard, exhaustive search, computer-assisted proof, OEIS A279405}
}

\newtheorem{theorem}{Theorem}
\newtheorem{lemma}[theorem]{Lemma}
\newtheorem{proposition}[theorem]{Proposition}
\newtheorem{corollary}[theorem]{Corollary}
\theoremstyle{remark}
\newtheorem{remark}[theorem]{Remark}

\newcommand{\Z}{\mathbb{Z}}
\newcommand{\Zn}{\Z/n\Z}
\newcommand{\Zf}{\Z/15\Z}
\newcommand{\Zs}{\Z/16\Z}
\newcommand{\Bl}{\mathrm{B}}
\newcommand{\Wh}{\mathrm{W}}

\title[Peaceable queens on the $15\times 15$ and $16\times 16$ tori]{Peaceable queens on the $15\times 15$ and $16\times 16$ tori:\\ the exact values $t(15)=20$ and $t(16)=32$}

\author{Jan Philipp Harries}
\thanks{ellamind GmbH, Bremen, Germany. \textit{E-mail address:} \href{mailto:jan@ellamind.com}{jan@ellamind.com}}

\date{July 25, 2026}

\subjclass[2020]{Primary 05B99; Secondary 05C69, 68R05, 05C30}
\keywords{peaceable queens, toroidal chessboard, exhaustive search, computer-assisted proof, OEIS A279405}

\raggedbottom

\begin{document}

\begin{abstract}
Let $t(n)$ denote the largest $m$ such that $m$ black and $m$ white queens can be placed on the $n\times n$ toroidal chessboard with no black queen attacking a white queen. We prove that $t(15)=20$ and $t(16)=32$, determining the first two previously unknown terms of this sequence (OEIS A279405), whose exact values were known only through $n=14$. Both lower bounds are explicit placements; the one for $n=16$ is the $(4,4)$-plaid of Clinch, Drescher, Huynh and Saffidine. Both upper bounds are computer-assisted exhaustive proofs on a common plan: a queen placement is replaced by a colouring of the $4n$ toroidal lines, elementary counting identities restrict the line-set cardinalities to finitely many \emph{size profiles} up to symmetry, and each profile is decided exactly by a finite search whose every pruning step is a necessary condition. For $n=15$ a profile is the quadruple of line-set sizes; there are $247$ of them modulo the induced dihedral action on the four line families and, in the top stratum, colour complementation --- an action induced by an affine board-symmetry group of order $14{,}400$ --- and the search examined $6{,}074{,}753{,}568$ cases. For $n=16$ the diagonal and antidiagonal sizes must be refined by parity, because on an even torus a diagonal and an antidiagonal meet in two cells or in none; there are $342$ canonical six-part profiles giving $677$ oriented search jobs, and the search examined $13{,}163{,}028{,}768$ cases. No larger placement was found in either search. The two values sharpen the parity contrast of the toroidal problem: $t(16)$ equals $16^2/8$ exactly, whereas $t(15)$ falls well short of $15^2/8$.
\end{abstract}

\maketitle

\section{Introduction}\label{sec:intro}

A \emph{peaceable} pair of queen armies is a placement of $m$ black and $m$ white queens on a chessboard such that no black queen attacks a white queen; queens of the same colour may attack each other freely. The problem of maximizing $m$ goes back to Ainley \cite{Ainley77} and was popularized through the OEIS sequence A250000 \cite{OEIS250000} for the ordinary $n\times n$ board. On the \emph{toroidal} board the four lines through a cell wrap around, which makes the problem homogeneous: the symmetry group acts transitively on cells, and the combinatorics becomes arithmetic in $\Zn$. The toroidal maxima form OEIS sequence A279405 \cite{OEIS279405}.

\begin{table}[h]
\centering
\begin{tabular}{@{}l*{9}{r}@{}}
\toprule
$n$      & 8 & 9 & 10 & 11 & 12 & 13 & 14 & \textbf{15} & \textbf{16}\\
$t(n)$   & 8 & 7 & 12 & 10 & 18 & 16 & 25 & \textbf{20} & \textbf{32}\\
$n^2/8$  & 8.0 & 10.1 & 12.5 & 15.1 & 18.0 & 21.1 & 24.5 & 28.1 & 32.0\\
\bottomrule
\end{tabular}
\medskip
\caption{The upper end of A279405. The values through $n=14$ are those recorded in \cite{OEIS279405}; $t(13)=16$ and $t(14)=25$ were contributed by Jaideep Padhi in June 2026. The values at $n=15$ and $n=16$ are Theorems~\ref{thm:main15} and~\ref{thm:main16}.}\label{tab:oeis}
\end{table}

Table~\ref{tab:oeis} shows the characteristic parity split of the toroidal problem: for even $n$ the known values sit at or near $n^2/8$ and at $n=14$ exceed it, whereas for odd $n$ they fall well short of it, hovering nearer $n^2/11$. Clinch, Drescher, Huynh and Saffidine \cite{CDHS24} made this contrast precise in the asymptotic regime, proving $t(n)\le 0.125\,n^2$ for odd $n$ \cite[Theorem~1.4]{CDHS24} against $t(n)\ge 0.1339\,n^2$ for all sufficiently large even $n$ \cite[Theorem~1.1]{CDHS24}, the latter by an explicit family of constructions they call \emph{plaids}. Their best published upper bound for the even torus is $t(n)\le 0.1402\,n^2$ for all even $n$ \cite[Theorem~1.2]{CDHS24}, the numerical optimum of the associated non-linear program being $0.140132\ldots$ \cite[\S4.2]{CDHS24}.

At $n=15$ and $n=16$ this left two gaps. For $n=15$, Theorem~1.4 of \cite{CDHS24} gives $t(15)\le 28$ by integrality, and the appendix of \cite{CDHS24} exhibits an explicit $20+20$ configuration, so $t(15)\in[20,28]$. For $n=16$, Theorem~1.2 of \cite{CDHS24} gives $t(16)\le 35$ by integrality, and the $(4,4)$-plaid on the $16\times16$ torus is peaceable with $32$ queens of each colour (Section~\ref{sec:lower16}), so $t(16)\in[32,35]$. A279405 listed no term for either $n$.

We determine both.

\begin{theorem}\label{thm:main15}
On the $15\times 15$ toroidal chessboard, $20$ black and $20$ white queens can be placed with no black--white attack, and $21$ black and $21$ white queens cannot. That is, $t(15)=20$.
\end{theorem}

\begin{theorem}\label{thm:main16}
On the $16\times 16$ toroidal chessboard, $32$ black and $32$ white queens can be placed with no black--white attack, and $33$ black and $33$ white queens cannot. That is, $t(16)=32$.
\end{theorem}

Both lower bounds are placements, reproduced in Sections~\ref{sec:lower15} and~\ref{sec:lower16}; neither is new. The upper bounds are the substance of this note and are computer-assisted. They follow a common plan, and the two cases differ in exactly the place where the arithmetic of $\Zn$ differs.

Section~\ref{sec:reduction} eliminates the queens, for every $n$ at once. A peaceable $m+m$ placement exists if and only if the $4n$ lines of the board can be coloured black and white so that the set of cells all four of whose lines are black, and the set of cells all four of whose lines are white, both have at least $m$ elements. The black line sets $R,C,D,A\subseteq\Zn$ (rows, columns, diagonals, antidiagonals) then determine everything, and swapping colours if necessary one may assume $S:=|R|+|C|+|D|+|A|\le 2n$. We stress that $S\le 2n$ is genuinely an inequality: it cannot be normalized to $S=2n$, because adding a line to the black side removes it from the white side and may destroy the white count. Section~\ref{sec:reduction} also records the two counting facts on which everything rests: which pairs of lines meet in how many cells, and an exact identity for triples of families.

The single arithmetic difference between the two cases is the incidence of a diagonal with an antidiagonal. The cells on diagonal $d$ and antidiagonal $a$ are the solutions of $2r=a+d$, $2c=a-d$. For odd $n$ the element $2$ is a unit, so there is exactly one such cell, and the four line families behave symmetrically. For even $n$ there are two such cells when $d\equiv a\pmod 2$ and none otherwise. This costs the even case one symmetry --- the halving map that exchanges $\{R,C\}$ with $\{D,A\}$ is no longer available --- and gains it a refinement: the diagonal and antidiagonal sizes must be split by the parity of the line label.

Accordingly, Sections~\ref{sec:lower15}--\ref{sec:result15} treat $n=15$ and Sections~\ref{sec:lower16}--\ref{sec:result16} treat $n=16$, in parallel. In each case one bounds the possible size profiles by counting (Sections~\ref{sec:profiles15},~\ref{sec:profiles16}), decides a single profile by an exact three-stage search in which every pruning rule is a necessary condition (Sections~\ref{sec:completion15},~\ref{sec:completion16}), and reports the outcome (Sections~\ref{sec:result15},~\ref{sec:result16}): $247$ canonical profiles and $6{,}074{,}753{,}568$ cases for $n=15$, all infeasible; $342$ canonical profiles, $677$ oriented jobs and $13{,}163{,}028{,}768$ cases for $n=16$, all infeasible.

Section~\ref{sec:verification} describes the verification, which deserves the most scrutiny and accordingly receives the most space, with a subsection for each case. The evidence is not equally strong on the two sides, and we say so there: for $n=15$ two independently written enumerators with different symmetry quotients and different completion algorithms agree on all $247$ profiles, whereas for $n=16$ the two programs that were run share everything except the completion kernel. Sections~\ref{sec:repro} and~\ref{sec:ai} give reproduction instructions and disclose the use of AI systems in producing this note.

\begin{remark}
Both proofs are exhaustive rather than explanatory. They give no human-readable reason why the odd torus is so much poorer than the even torus, and they do not sharpen the general bounds of \cite{CDHS24}. A structural proof of $t(15)\le 20$ --- or, better, of a general odd-$n$ bound near $n^2/11$ --- remains open and would supersede the first computation; likewise a structural proof that the plaid construction is optimal on the even torus would supersede the second.
\end{remark}

\begin{remark}\label{rem:conj51}
Conjecture~5.1 of \cite{CDHS24} predicts $t(n)\le 2+\max\min\{|B|,|W|\}$ for even $n$, the maximum taken over $(a,b)$-plaids on the $n\times n$ torus. A direct check of the $64$ plaids available at $n=16$ gives $\max\min\{|B|,|W|\}=32$, so the conjecture predicts $t(16)\le 34$. Theorem~\ref{thm:main16} confirms the conjecture at $n=16$ and shows that its slack of $2$ is not attained there.
\end{remark}

\section{Notation and the line-colouring reduction}\label{sec:reduction}

Throughout, $n\ge 2$ and all coordinates lie in $\Zn$, represented by $\{0,1,\dots,n-1\}$. The board is $G=(\Zn)^2$, with $n^2$ cells. Through a cell $(r,c)$ pass four \emph{lines}:
\[
\text{row } r,\qquad \text{column } c,\qquad \text{diagonal } d=r-c,\qquad \text{antidiagonal } a=r+c ,
\]
each family indexed by $\Zn$, so the board carries $4n$ lines, each containing exactly $n$ cells. Two queens \emph{attack} each other exactly when they share a line. We write $t(n)$ for the largest $m$ such that $m$ black and $m$ white queens can be placed on the $n\times n$ torus with no black--white attack. All $2m$ queens occupy distinct cells; queens of one colour may share lines.

For $R,C,D,A\subseteq\Zn$ put
\[
\Bl(R,C,D,A)=\{(r,c)\in G: r\in R,\ c\in C,\ r-c\in D,\ r+c\in A\},
\]
and let $\Wh(R,C,D,A)=\Bl(R^{c},C^{c},D^{c},A^{c})$, where ${}^{c}$ denotes complement in $\Zn$ taken family by family. Cells of $\Bl$ are called \emph{black-homogeneous} and cells of $\Wh$ \emph{white-homogeneous}.

\begin{lemma}[Line-colouring equivalence]\label{lem:lines}
Let $m\ge 1$. There is a peaceable placement of $m$ black and $m$ white queens on the $n\times n$ torus if and only if there are sets $R,C,D,A\subseteq\Zn$ with
\[
|\Bl(R,C,D,A)|\ge m
\quad\text{and}\quad
|\Wh(R,C,D,A)|\ge m .
\]
\end{lemma}

\begin{proof}
($\Rightarrow$) Given a peaceable placement, colour every line carrying a black queen black and every line carrying a white queen white. No line receives both colours: a line carrying queens of both colours would exhibit a black--white attack. Colour every remaining line white, and let $R,C,D,A$ be the black line sets in the four families; the white line sets are then exactly the complements $R^{c},C^{c},D^{c},A^{c}$. Every black queen sits on four black lines, so its cell lies in $\Bl$, and the $m$ black queens occupy distinct cells, whence $|\Bl|\ge m$; likewise every white queen sits on four lines that carry no black queen, hence on four white lines, so $|\Wh|\ge m$.

($\Leftarrow$) Given such sets, place black queens on any $m$ cells of $\Bl$ and white queens on any $m$ cells of $\Wh$. The two sets are disjoint, since a cell cannot have all four of its lines both in and out of the black sets. If a black queen attacked a white queen they would share a line, say a row $r$; but then $r\in R$ and $r\in R^{c}$, which is absurd.
\end{proof}

Colouring the unused lines white rather than leaving them unused is harmless in the direction that matters: it can only enlarge the white sets, and $\Wh$ is monotone increasing in the complements. From now on, a \emph{configuration} is a quadruple $(R,C,D,A)$ of subsets of $\Zn$. Write
\[
s_0=|R|,\quad s_1=|C|,\quad s_2=|D|,\quad s_3=|A|,\qquad S=s_0+s_1+s_2+s_3 .
\]

\begin{lemma}[Budget]\label{lem:budget}
If some configuration has $|\Bl|\ge m$ and $|\Wh|\ge m$, then some configuration with the same property has $S\le 2n$.
\end{lemma}

\begin{proof}
Replacing $(R,C,D,A)$ by $(R^{c},C^{c},D^{c},A^{c})$ exchanges $\Bl$ and $\Wh$, hence preserves the property, and replaces $S$ by $4n-S$. One of $S$ and $4n-S$ is at most $2n$.
\end{proof}

\begin{remark}\label{rem:notequal}
Lemma~\ref{lem:budget} is used only as an inequality. We do not normalize to $S=2n$: no lemma asserting that a feasible configuration can always be modified into one with $S=2n$ is proved or used here, and none is obviously true, because feasibility is not monotone under transferring a line from the white side to the black side --- a line added to $R$ is removed from $R^{c}$ and can destroy $|\Wh|\ge m$. (It can also fail to: the witness of Proposition~\ref{prop:lower15}, with black-side $S=28$, happens to extend to a configuration with $S=30$ and $|\Bl|=|\Wh|=20$ by adding row $2$ and column $11$, so individual examples decide nothing either way.) Accordingly both searches below conservatively cover the whole admissible range of $S$, not merely the top stratum. For $n=16$ the necessary conditions of Section~\ref{sec:profiles16} turn out to admit no profile at all with $S\in\{30,31,32\}$.
\end{remark}

\begin{lemma}[Incidence]\label{lem:incidence}
Let $X$ and $Y$ be two distinct line families and let $\ell$ be a line of $X$ and $\ell'$ a line of $Y$.
\begin{enumerate}
\item[(i)] If $\{X,Y\}\ne\{D\text{-family},A\text{-family}\}$, then $\ell$ and $\ell'$ meet in exactly one cell.
\item[(ii)] If $\ell$ is the diagonal $d$ and $\ell'$ the antidiagonal $a$, their common cells are the solutions of
\begin{equation}\label{eq:da}
2r=a+d,\qquad 2c=a-d \qquad\text{in }\Zn .
\end{equation}
For odd $n$ there is exactly one. For even $n$ there are exactly two if $d\equiv a\pmod 2$, and none otherwise; the two are $(r,c)$ and $(r+\tfrac n2,\,c+\tfrac n2)$.
\end{enumerate}
\end{lemma}

\begin{proof}
(i) A row $r$ and a column $c$ meet in $(r,c)$; a row $r$ and a diagonal $d$ in $(r,r-d)$; a row $r$ and an antidiagonal $a$ in $(r,a-r)$; a column $c$ and a diagonal $d$ in $(c+d,c)$; a column $c$ and an antidiagonal $a$ in $(a-c,c)$. In each case the cell is unique.

(ii) A cell $(r,c)$ lies on diagonal $d$ and antidiagonal $a$ iff $r-c=d$ and $r+c=a$; adding and subtracting gives \eqref{eq:da}, and conversely \eqref{eq:da} together with either original equation gives the other. Consider the homomorphism $\varphi\colon G\to(\Zn)^2$, $\varphi(r,c)=(r-c,r+c)$. Its kernel is $\{(r,c):r=c,\ 2r=0\}$. For odd $n$ the kernel is trivial and $\varphi$ is a bijection, with inverse $r=2^{-1}(d+a)$, $c=2^{-1}(a-d)$; so each fibre is a single cell. For even $n$ the kernel is $\{(0,0),(\tfrac n2,\tfrac n2)\}$, of order $2$, so every nonempty fibre has exactly two elements and the image has $n^2/2$ elements. Every $(d,a)$ in the image satisfies $d+a=2r\equiv 0\pmod 2$, and the set of such pairs also has $n^2/2$ elements; hence the image is exactly $\{(d,a):d\equiv a\pmod2\}$. The description of the second solution is the kernel translate.
\end{proof}

For two families $X,Y$ with black sets $\mathcal X,\mathcal Y$ write $m(X,Y)$ for the number of cells lying on a black line of $X$ and on a black line of $Y$. By Lemma~\ref{lem:incidence}, $m(X,Y)=|\mathcal X|\,|\mathcal Y|$ unless $\{X,Y\}$ is the diagonal--antidiagonal pair; the exceptional value is computed in each case below.

\begin{lemma}[Triple identity]\label{lem:triples}
Let $X,Y,Z$ be three of the four families, with black sets of sizes $x,y,z$. Let $T$ be the number of cells whose lines in these three families are all black, and $T'$ the number whose lines in these three families are all white. Then
\begin{equation}\label{eq:triple-id}
T+T'=n^2-n(x+y+z)+m(X,Y)+m(X,Z)+m(Y,Z) .
\end{equation}
Consequently, if $|\Bl|\ge m$ and $|\Wh|\ge m$ then the right-hand side of \eqref{eq:triple-id} is at least $2m$, for each of the four triples of families.
\end{lemma}

\begin{proof}
Write $u,v,w\colon G\to\{0,1\}$ for the indicators of ``the cell's line in family $X$ (resp.\ $Y$, $Z$) is black''. Then
\[
uvw+(1-u)(1-v)(1-w)=1-u-v-w+uv+uw+vw .
\]
Summing over the $n^2$ cells: $\sum 1=n^2$; $\sum u=nx$, since each of the $x$ black lines of family $X$ contains $n$ cells, and likewise $\sum v=ny$, $\sum w=nz$; and $\sum uv=m(X,Y)$ by definition, likewise for the other two pairs. This is \eqref{eq:triple-id}. Every cell of $\Bl$ is counted by $T$ and every cell of $\Wh$ by $T'$, and the two sets are disjoint, so $|\Bl|+|\Wh|\le T+T'$.
\end{proof}

\medskip
\noindent\textbf{Part I: the odd case $n=15$.}

\section{The lower bound for $n=15$}\label{sec:lower15}

\begin{proposition}\label{prop:lower15}
$t(15)\ge 20$.
\end{proposition}

\begin{proof}
Take the black line sets
\begin{equation}\label{eq:witness15}
\begin{aligned}
R&=\{5,7,8,11,12,14\}, & D&=\{0,1,4,8,12,13,14\},\\
C&=\{0,3,4,7,9,10,12,13\}, & A&=\{1,2,3,9,10,12,14\},
\end{aligned}
\end{equation}
of sizes $(6,8,7,7)$ in $\Zf$, and colour every remaining line white, giving white line sets of sizes $(9,7,8,8)$. Direct enumeration of the $225$ cells gives exactly $20$ black-homogeneous cells,
\[
\begin{array}{llll}
(5,4)&(5,7)&(5,12)&\\
(7,3)&(7,7)&(7,9)&(7,10)\\
(8,4)&(8,9)&(8,10)&\\
(11,3)&(11,7)&(11,13)&\\
(12,0)&(12,4)&(12,12)&(12,13)\\
(14,0)&(14,10)&(14,13)&
\end{array}
\]
and exactly $22$ white-homogeneous cells, of which we use the following $20$:
\[
\begin{array}{llll}
(0,5)&(0,6)&(0,8)&\\
(1,5)&(1,6)&(1,14)&\\
(2,6)&(2,11)&&\\
(3,1)&(3,8)&&\\
(4,1)&(4,2)&&\\
(6,1)&(6,14)&&\\
(9,2)&(9,6)&(9,14)&\\
(10,1)&(10,5)&(13,2)&
\end{array}
\]
(the two unused ones are $(13,6)$ and $(13,8)$). Place black queens on the first list and white queens on the second. Every black queen lies on four black lines and every white queen on four white lines; a black--white attack would require a line that is both black and white, and there is none. Hence $20$ black and $20$ white queens coexist peaceably.
\end{proof}

Figure~\ref{fig:board15} displays the configuration. This lower bound is not new: an explicit $20+20$ configuration for $n=15$ is given in the appendix of \cite{CDHS24}. We reproduce a placement here in line-set form so that Theorem~\ref{thm:main15} is self-contained and checkable by hand.

\begin{figure}[t]
\centering
\setlength{\tabcolsep}{2.4pt}
\renewcommand{\arraystretch}{1.05}
{\small
\begin{tabular}{r|*{15}{c}}
$r\backslash c$ & 0 & 1 & 2 & 3 & 4 & 5 & 6 & 7 & 8 & 9 & 10 & 11 & 12 & 13 & 14 \\
\midrule
0 & $\cdot$ & $\cdot$ & $\cdot$ & $\cdot$ & $\cdot$ & $\square$ & $\square$ & $\cdot$ & $\square$ & $\cdot$ & $\cdot$ & $\cdot$ & $\cdot$ & $\cdot$ & $\cdot$ \\
1 & $\cdot$ & $\cdot$ & $\cdot$ & $\cdot$ & $\cdot$ & $\square$ & $\square$ & $\cdot$ & $\cdot$ & $\cdot$ & $\cdot$ & $\cdot$ & $\cdot$ & $\cdot$ & $\square$ \\
2 & $\cdot$ & $\cdot$ & $\cdot$ & $\cdot$ & $\cdot$ & $\cdot$ & $\square$ & $\cdot$ & $\cdot$ & $\cdot$ & $\cdot$ & $\square$ & $\cdot$ & $\cdot$ & $\cdot$ \\
3 & $\cdot$ & $\square$ & $\cdot$ & $\cdot$ & $\cdot$ & $\cdot$ & $\cdot$ & $\cdot$ & $\square$ & $\cdot$ & $\cdot$ & $\cdot$ & $\cdot$ & $\cdot$ & $\cdot$ \\
4 & $\cdot$ & $\square$ & $\square$ & $\cdot$ & $\cdot$ & $\cdot$ & $\cdot$ & $\cdot$ & $\cdot$ & $\cdot$ & $\cdot$ & $\cdot$ & $\cdot$ & $\cdot$ & $\cdot$ \\
5 & $\cdot$ & $\cdot$ & $\cdot$ & $\cdot$ & $\blacksquare$ & $\cdot$ & $\cdot$ & $\blacksquare$ & $\cdot$ & $\cdot$ & $\cdot$ & $\cdot$ & $\blacksquare$ & $\cdot$ & $\cdot$ \\
6 & $\cdot$ & $\square$ & $\cdot$ & $\cdot$ & $\cdot$ & $\cdot$ & $\cdot$ & $\cdot$ & $\cdot$ & $\cdot$ & $\cdot$ & $\cdot$ & $\cdot$ & $\cdot$ & $\square$ \\
7 & $\cdot$ & $\cdot$ & $\cdot$ & $\blacksquare$ & $\cdot$ & $\cdot$ & $\cdot$ & $\blacksquare$ & $\cdot$ & $\blacksquare$ & $\blacksquare$ & $\cdot$ & $\cdot$ & $\cdot$ & $\cdot$ \\
8 & $\cdot$ & $\cdot$ & $\cdot$ & $\cdot$ & $\blacksquare$ & $\cdot$ & $\cdot$ & $\cdot$ & $\cdot$ & $\blacksquare$ & $\blacksquare$ & $\cdot$ & $\cdot$ & $\cdot$ & $\cdot$ \\
9 & $\cdot$ & $\cdot$ & $\square$ & $\cdot$ & $\cdot$ & $\cdot$ & $\square$ & $\cdot$ & $\cdot$ & $\cdot$ & $\cdot$ & $\cdot$ & $\cdot$ & $\cdot$ & $\square$ \\
10 & $\cdot$ & $\square$ & $\cdot$ & $\cdot$ & $\cdot$ & $\square$ & $\cdot$ & $\cdot$ & $\cdot$ & $\cdot$ & $\cdot$ & $\cdot$ & $\cdot$ & $\cdot$ & $\cdot$ \\
11 & $\cdot$ & $\cdot$ & $\cdot$ & $\blacksquare$ & $\cdot$ & $\cdot$ & $\cdot$ & $\blacksquare$ & $\cdot$ & $\cdot$ & $\cdot$ & $\cdot$ & $\cdot$ & $\blacksquare$ & $\cdot$ \\
12 & $\blacksquare$ & $\cdot$ & $\cdot$ & $\cdot$ & $\blacksquare$ & $\cdot$ & $\cdot$ & $\cdot$ & $\cdot$ & $\cdot$ & $\cdot$ & $\cdot$ & $\blacksquare$ & $\blacksquare$ & $\cdot$ \\
13 & $\cdot$ & $\cdot$ & $\square$ & $\cdot$ & $\cdot$ & $\cdot$ & $\lozenge$ & $\cdot$ & $\lozenge$ & $\cdot$ & $\cdot$ & $\cdot$ & $\cdot$ & $\cdot$ & $\cdot$ \\
14 & $\blacksquare$ & $\cdot$ & $\cdot$ & $\cdot$ & $\cdot$ & $\cdot$ & $\cdot$ & $\cdot$ & $\cdot$ & $\cdot$ & $\blacksquare$ & $\cdot$ & $\cdot$ & $\blacksquare$ & $\cdot$ \\
\end{tabular}}
\caption{The $20+20$ placement of Proposition~\ref{prop:lower15}: $\blacksquare$ black queens, $\square$ white queens, $\lozenge$ the two white-homogeneous cells left empty. Rows, columns and both diagonal families wrap around.}\label{fig:board15}
\end{figure}

\section{Admissible size profiles for $n=15$}\label{sec:profiles15}

Here $n=15$, $2$ is a unit modulo $15$ with $2^{-1}=8$, and by Lemma~\ref{lem:incidence} any two lines from distinct families meet in exactly one cell; in particular $m(X,Y)=|\mathcal X||\mathcal Y|$ for all six pairs. We call $(s_0,s_1,s_2,s_3)$ the \emph{size profile} of the configuration and ask whether a configuration exists with $|\Bl|\ge 21$ and $|\Wh|\ge 21$.

\begin{lemma}[Pair bounds]\label{lem:pairs15}
Let $(R,C,D,A)$ be a configuration in $\Zf$ with size profile $(s_0,s_1,s_2,s_3)$. Then for all $i<j$,
\[
|\Bl|\le s_is_j,\qquad |\Wh|\le (15-s_i)(15-s_j).
\]
Consequently, if $|\Bl|\ge 21$ and $|\Wh|\ge 21$ then
\begin{equation}\label{eq:pair15}
s_is_j\ge 21\quad\text{and}\quad (15-s_i)(15-s_j)\ge 21\qquad\text{for all }i<j .
\end{equation}
\end{lemma}

\begin{proof}
Every cell of $\Bl$ lies on a black line of family $i$ and on a black line of family $j$, and distinct cells of $\Bl$ give distinct such pairs of lines, because by Lemma~\ref{lem:incidence}(i)--(ii) two lines from distinct families meet in exactly one cell when $n$ is odd; hence $|\Bl|\le s_is_j$. The bound for $\Wh$ is the same statement applied to the complementary sets.
\end{proof}

By Lemma~\ref{lem:triples} with $m(X,Y)=xy$ throughout, if $|\Bl|\ge 21$ and $|\Wh|\ge 21$ then
\begin{equation}\label{eq:triple15}
225-15(x+y+z)+xy+xz+yz\ge 42
\end{equation}
for each of the four triples of families, where $x,y,z$ are the corresponding sizes. Note that \eqref{eq:triple15} is symmetric under $(x,y,z)\mapsto(15-x,15-y,15-z)$, as it must be.

\subsection*{Symmetry}

Let $\Gamma$ be the group of affine bijections of $G$ that permute the four line families. It consists of the $225$ translations composed with $64$ linear maps, so $|\Gamma|=14{,}400$; this was verified by exhaustive enumeration of the invertible $2\times2$ matrices over $\Zf$. Every $\gamma\in\Gamma$ carries a configuration $(R,C,D,A)$ to a configuration with the same $|\Bl|$ and $|\Wh|$, permuting the families according to a homomorphism $\Gamma\to\mathrm{Sym}\{R,C,D,A\}$. The image is the dihedral group of order $8$ preserving the partition $\{\{R,C\},\{D,A\}\}$, generated by
\begin{align*}
\text{transpose } \tau(r,c)&=(c,r): && R\leftrightarrow C,\ D\mapsto -D,\ A\mapsto A;\\
\text{reflection } \nu(r,c)&=(r,-c): && R\mapsto R,\ C\mapsto -C,\ D\leftrightarrow A;\\
\text{halving } \psi(r,c)&=(r+c,\ r-c): && \{R,C\}\leftrightarrow\{D,A\} .
\end{align*}
For $\psi$ one checks directly that a level set of the new row coordinate is an old antidiagonal, of the new column coordinate an old diagonal, of the new diagonal an old column and of the new antidiagonal an old row; $\psi$ is bijective because $\det\begin{pmatrix}1&1\\1&-1\end{pmatrix}=-2$ is a unit modulo $15$. Composing $\tau$, $\nu$ and $\psi$ yields, in one-line notation listing the images of $(R,C,D,A)$, exactly the eight family permutations
\[
RCDA,\quad RCAD,\quad CRDA,\quad CRAD,\quad DARC,\quad DACR,\quad ADRC,\quad ADCR,
\]
which one verifies by direct substitution to be closed under composition: they form the dihedral group of order $8$ preserving the partition $\{\{R,C\},\{D,A\}\}$, and no composition of the three generators produces a permutation splitting that partition. Conversely, every family-preserving linear map agrees with one of eight representatives --- one above each induced permutation --- up to composition with a common unit scaling $(r,c)\mapsto(ur,uc)$, and $|(\Zf)^{\times}|=8$, which accounts for all $8\cdot 8=64$ linear maps found by the exhaustive matrix enumeration. The induced action on size profiles is the permutation action of this $D_4$ on ordered quadruples.

\begin{proposition}\label{prop:247}
Let $(R,C,D,A)$ be a configuration in $\Zf$ with $|\Bl|\ge 21$, $|\Wh|\ge 21$ and $S\le 30$. Then its size profile is equivalent, under the $D_4$ action together with complementation $s_i\mapsto 15-s_i$ in the case $S=30$, to one of exactly $247$ canonical profiles, distributed by $S$ as in Table~\ref{tab:strata15}.
\end{proposition}

\begin{table}[h]
\centering
\begin{tabular}{@{}l*{11}{r}@{}}
\toprule
$S$ & 20 & 21 & 22 & 23 & 24 & 25 & 26 & 27 & 28 & 29 & 30 \\
canonical profiles & 1 & 1 & 4 & 6 & 13 & 18 & 29 & 37 & 47 & 52 & 39 \\
\bottomrule
\end{tabular}
\medskip
\caption{The $247$ canonical size profiles for $n=15$, by budget $S$.}\label{tab:strata15}
\end{table}

\begin{proof}
Filter all $16^4$ quadruples $(s_0,s_1,s_2,s_3)\in\{0,\dots,15\}^4$ by $S\le 30$ and by conditions \eqref{eq:pair15} and \eqref{eq:triple15}, which are necessary by Lemmas~\ref{lem:pairs15} and~\ref{lem:triples}, and canonicalize the survivors: replace each by the lexicographically least element of its $D_4$ orbit, adjoining the $D_4$ orbit of the complementary profile $(15-s_i)_i$ when $S=30$. Complementation preserves the property ``$|\Bl|\ge21$ and $|\Wh|\ge21$'' but sends $S$ to $60-S$, so it may be used as an identification only within the stratum $S=30$, which it preserves. The enumeration is a finite computation over $65{,}536$ quadruples; it was carried out independently three times, in two languages, always giving the same $247$ profiles and the distribution of Table~\ref{tab:strata15}.
\end{proof}

Inspection of the list shows that every surviving profile satisfies $3\le s_i\le 12$ for all $i$; the unique profiles at the two smallest budgets are $(5,5,5,5)$ and $(5,5,5,6)$, while the top stratum $S=30$ contains both the lopsided $(3,8,9,10)$ and the balanced $(7,7,8,8)$.

\section{Deciding one profile exactly for $n=15$}\label{sec:completion15}

Fix a canonical profile. By the $D_4$ symmetry we may choose either $(R,C)$ or $(D,A)$ as the \emph{fixed pair}; write $(r,c)$ for its two sizes, with $r\le c$ after a transpose if necessary, and $(d,a)$ for the sizes of the two remaining families, which we call the \emph{completion} families. The search proceeds in three nested stages: orbit representatives for the fixed pair, all subsets for one completion family, and an exact decision for the last.

\subsection*{Stage 1: representatives of the fixed pair}

Fix the pair $(R,C)$ up to a subgroup $H\le\Gamma$ of the stabilizer of the ordered profile that also preserves the chosen fixed-pair/completion decomposition. It is generated by:
\begin{itemize}
\item independent translations $(R,C)\mapsto(R+\alpha,\ C+\beta)$, which change $D$ and $A$ only by translations ($225$ elements);
\item common unit scaling $(R,C)\mapsto(uR,\ uC)$ for $u\in(\Zf)^{\times}$, which scales $D$ and $A$ by the same $u$ ($8$ elements);
\item the transpose $\tau$, admissible exactly when $r=c$, since $\tau$ exchanges the two fixed families while sending $D\mapsto -D$ and fixing $A$, hence preserving the completion sizes;
\item the relative sign $\nu$, i.e.\ $(R,C)\mapsto(R,-C)$, admissible \emph{only} when $d=a$.
\end{itemize}
The last restriction is the one soundness point that is easy to get wrong. Since $\nu$ exchanges the diagonal and antidiagonal families, an ordered profile with $d\ne a$ is not mapped to itself, and quotienting by $\nu$ would then identify configurations belonging to different ordered profiles --- legitimate only if both orientations are enumerated. We use $\nu$ only when $d=a$, where it is a genuine automorphism of the ordered profile. An earlier cost model in this project applied $\nu$ unconditionally; correcting it raised the workload of the independent enumerator from $4{,}876{,}723{,}839$ to $6{,}241{,}793{,}402$ completions.

$H$ need not be the full stabilizer of the ordered profile inside $\Gamma$: when all four sizes are equal, for instance, $\psi$ stabilizes the ordered profile but exchanges the fixed pair with the completion pair, and it is deliberately not in $H$. Quotienting by a subgroup is always sound --- it can only produce more representatives than a larger group would, never miss an orbit --- so omitted symmetries cost work, not correctness. The quotient is justified by the following transfer lemma.

\begin{lemma}[Orbit transfer]\label{lem:transfer15}
Let $(R,C)$ and $(R',C')$ lie in the same $H$-orbit. Then for each pair of prescribed completion sizes $(d,a)$, a completion $(D,A)$ of $(R,C)$ with $|\Bl|\ge 21$, $|\Wh|\ge 21$ and $|D|=d$, $|A|=a$ exists if and only if the corresponding completion of $(R',C')$ exists.
\end{lemma}

\begin{proof}
An element of $H$ with $(R,C)=h(R',C')$ is induced by an affine board symmetry $\gamma\in\Gamma$. If $(R',C',D',A')$ is a configuration, then $\gamma$ carries it to a configuration with fixed pair $(R,C)$ and the same $|\Bl|$ and $|\Wh|$, and by the admissibility conditions above ($\tau$ only when $r=c$, $\nu$ only when $d=a$) the completion sizes are preserved. Apply this to a feasible completion in either direction.
\end{proof}

By Lemma~\ref{lem:transfer15} it suffices that the enumerated set of pairs intersect every $H$-orbit of pairs of the prescribed sizes; nothing more is required for soundness. The group has order $225\cdot 8=1800$, doubled for each of the two admissible extra generators, hence $1800$, $3600$ or $7200$. The implementation produces its pair set by reducing each coordinate to its cyclic-translation necklace and then sweeping all necklace pairs: each pair not previously marked is emitted as a representative, and all images of an emitted pair under the residual unit, transpose and sign actions are marked. A pair is skipped only when marked, i.e.\ only when an explicit $h\in H$ mapping an emitted representative to it has been exhibited, so the emitted set meets every orbit. That it is moreover an exact transversal --- one representative per orbit --- is not needed for soundness but is verified anyway: only $43$ distinct triples $(r,c,\text{group})$ occur across the $247$ profiles, and for each of them the cardinality of the emitted set equals the orbit count computed independently by Burnside's lemma (Section~\ref{sec:verification}).

Choosing the fixed pair is a free optimization: for each profile both choices are costed as (number of pair representatives) $\times$ (number of subsets of the smaller completion family) and the cheaper one is taken. Likewise, within the completion families, the family with fewer subsets is the one enumerated in Stage~2, which is legitimate because $\nu$ exchanges the two completion families globally, so the ordered profiles $(r,c,d,a)$ and $(r,c,a,d)$ are simultaneously feasible or infeasible.

\subsection*{Stage 2: the completion as a subset problem}

Fix a representative $(R,C)$. Index cells by their diagonal and antidiagonal coordinates via the inverse $r=8(d+a)$, $c=8(a-d)$ of Lemma~\ref{lem:incidence}(ii), writing $x$ for the diagonal and $y$ for the antidiagonal coordinate to keep them apart from the sizes $d,a$, and set
\[
P_{x,y}=\mathbf 1_R\bigl(8(x+y)\bigr)\,\mathbf 1_C\bigl(8(y-x)\bigr),
\qquad
Q_{x,y}=\mathbf 1_{R^{c}}\bigl(8(x+y)\bigr)\,\mathbf 1_{C^{c}}\bigl(8(y-x)\bigr)
\qquad (x,y\in\Zf).
\]
For each antidiagonal set $A\subseteq\Zf$ of the prescribed size $a$, put
\[
p_x=\sum_{y\in A}P_{x,y},
\qquad
q_x=\sum_{y\notin A}Q_{x,y},
\qquad
Q_0=\sum_{x\in\Zf}q_x .
\]
Then for any diagonal set $D$,
\[
|\Bl(R,C,D,A)|=\sum_{x\in D}p_x,
\qquad
|\Wh(R,C,D,A)|=\sum_{x\notin D}q_x=Q_0-\sum_{x\in D}q_x .
\]
Hence a diagonal set $D$ of the prescribed size $d$ completes the configuration if and only if
\begin{equation}\label{eq:subset15}
\sum_{x\in D}p_x\ \ge\ 21
\qquad\text{and}\qquad
\sum_{x\in D}q_x\ \le\ Q_0-21 .
\end{equation}
This is an exact $15$-item, fixed-cardinality, two-constraint subset problem, with nonnegative integer weights $p_x,q_x\le 15$. Before it is posed, two scalar necessary conditions --- $\sum_x p_x\ge 21$ and $Q_0\ge 21$, both depending on $A$ alone --- are tested and cheaply discard most candidates $A$.

\subsection*{Stage 3: exact completion decision}

Problem \eqref{eq:subset15} is decided in two layers, both of which reject only by necessary conditions. First, two cardinality-aware root bounds: let $\mu$ be the sum of the $d$ largest of the weights $p_x$ and let $\lambda$ be the sum of the $d$ smallest of the weights $q_x$. Every $d$-element set $D$ satisfies $\sum_{x\in D}p_x\le\mu$ and $\sum_{x\in D}q_x\ge\lambda$, so if $\mu<21$ or $\lambda>Q_0-21$ the instance is infeasible and is rejected without further work; this layer is where the bulk of the scalar-prefilter survivors die (Table~\ref{tab:stats15}). Second, for the instances that pass both root bounds, a depth-first search over the $15$ items, each included or excluded, with the following bounds precomputed by suffix dynamic programming: for every position $x$ and every remaining cardinality $k$, the largest attainable sum of $k$ of the weights $p_x,\dots,p_{14}$, and the smallest attainable sum of $k$ of the weights $q_x,\dots,q_{14}$.

\begin{lemma}\label{lem:exhaustive15}
The two-layer decision of Stage~3 answers \eqref{eq:subset15} correctly.
\end{lemma}

\begin{proof}
The root bounds reject only infeasible instances, as just argued. In the search, a node is abandoned only when one of the following holds: the number of remaining items is smaller than the cardinality still required; the accumulated $q$-sum already exceeds $Q_0-21$; the accumulated $p$-sum plus the maximal attainable remaining $p$-sum with the required cardinality is below $21$; or the accumulated $q$-sum plus the minimal attainable remaining $q$-sum with the required cardinality exceeds $Q_0-21$. Each is a necessary condition for a feasible completion within the subtree, so no feasible $D$ is discarded. Conversely the search reports success only at a node where the required cardinality is reached and the $p$-constraint is met, or at a node where the $p$-constraint is already met and the minimal remaining $q$-sum with the required cardinality keeps the $q$-constraint satisfiable; in the latter case an explicit feasible $D$ exists, since $p$-sums only increase as further items are added. In both cases a feasible $D$ exists. As the search is otherwise a complete binary recursion over the $15$ items, it is exhaustive.
\end{proof}

Item order affects only speed, not correctness; the implementation sorts items by a ratio heuristic. Consequently the three stages together decide, for a given profile, whether \emph{any} configuration with that ordered size profile has $|\Bl|\ge 21$ and $|\Wh|\ge 21$, and the answer transfers to the whole $D_4$ orbit and, in the stratum $S=30$, to the complementary profile as well.

\section{The search and its outcome for $n=15$}\label{sec:result15}

The following lemma assembles the chain of Sections~\ref{sec:reduction}--\ref{sec:completion15} into a single soundness statement for the computation.

\begin{lemma}[End-to-end exhaustiveness, $n=15$]\label{lem:endtoend15}
If a configuration in $\Zf$ with $|\Bl|\ge 21$ and $|\Wh|\ge 21$ exists, then the search of Section~\ref{sec:completion15}, executed on all $247$ canonical profiles, reports a feasible completion for at least one of them.
\end{lemma}

\begin{proof}
By Lemma~\ref{lem:budget} such a configuration exists with $S\le 30$. Its size profile satisfies the pair and triple conditions (Lemmas~\ref{lem:pairs15} and~\ref{lem:triples}), so by Proposition~\ref{prop:247} some element $\gamma$ of the $D_4$ action --- composed, possibly, with complementation when $S=30$ --- carries the configuration to one whose ordered profile is a canonical worklist profile; $\gamma$ preserves feasibility. Within that profile's search, a further symmetry (the choice of fixed pair, Lemma~\ref{lem:transfer15}, and if applicable the Stage-2 relabelling of the completion families) carries the configuration to one whose fixed pair is an enumerated representative. Its Stage-2 set is among the exhaustively enumerated subsets of the prescribed size, and it passes the scalar conditions of Stage~2, which are necessary. Its remaining line set satisfies \eqref{eq:subset15}, so the instance passes the root bounds and, by Lemma~\ref{lem:exhaustive15}, the Stage-3 decision reports it feasible.
\end{proof}

\begin{proposition}\label{prop:unsat15}
For every one of the $247$ canonical profiles of Proposition~\ref{prop:247}, no configuration in $\Zf$ with that profile has $|\Bl|\ge 21$ and $|\Wh|\ge 21$.
\end{proposition}

\begin{proof}
The search of Section~\ref{sec:completion15} was executed on all $247$ profiles by \texttt{peace15\_solver.cpp} (SHA-256 digest in Section~\ref{sec:repro}) and reported no feasible completion anywhere; by Lemma~\ref{lem:endtoend15} this establishes the claim. The run statistics are given in Table~\ref{tab:stats15}; a per-profile certificate recording, for each profile, the fixed pair, the orbit count, a hash of the representative list, and the case counts, is included among the ancillary files, and the run was independently repeated as described in Section~\ref{sec:ver15}.
\end{proof}

\begin{table}[h]
\centering
\begin{tabular}{@{}lr@{}}
\toprule
canonical profiles                 & $247$ \\
distinct fixed-pair orbit types    & $43$ \\
Stage-2 completion subsets examined & $6{,}074{,}753{,}568$ \\
survivors of the scalar prefilter of Stage~2 & $2{,}329{,}772{,}941$ \\
survivors of the root bounds of Stage~3 (searches entered) & $26{,}948{,}453$ \\
branch nodes in those searches     & $246{,}476{,}279$ \\
profiles admitting a $21+21$ configuration & $0$ \\
\bottomrule
\end{tabular}
\medskip
\caption{Statistics of the exhaustive run for $n=15$, one line per pruning layer: of the $6{,}074{,}753{,}568$ enumerated Stage-2 subsets, $2{,}329{,}772{,}941$ pass the two scalar conditions of Stage~2, of which $26{,}948{,}453$ pass the two cardinality-aware root bounds of Stage~3 and enter the depth-first search; every layer rejects by necessary conditions only. The number of Stage-2 subsets examined equals $\sum_{\text{profiles}} (\text{orbit count})\cdot\binom{15}{a}$, which is checked row by row by the audit script.}\label{tab:stats15}
\end{table}

\begin{proof}[Proof of Theorem~\ref{thm:main15}]
The lower bound is Proposition~\ref{prop:lower15}. For the upper bound, suppose $21$ black and $21$ white queens are placed peaceably. By Lemma~\ref{lem:lines} there is a configuration with $|\Bl|\ge 21$ and $|\Wh|\ge 21$, by Lemma~\ref{lem:budget} one may take $S\le 30$, and by Proposition~\ref{prop:247} its size profile lies in the $247$-element canonical list. This contradicts Proposition~\ref{prop:unsat15}. Hence $t(15)\le 20$, and $t(15)=20$.
\end{proof}

\medskip
\noindent\textbf{Part II: the even case $n=16$.}

\section{The lower bound for $n=16$}\label{sec:lower16}

\begin{proposition}\label{prop:lower16}
$t(16)\ge 32$.
\end{proposition}

\begin{proof}
Work in $\Zs$ and take the black line sets
\begin{equation}\label{eq:witness16}
R=C=\{5,7,9,11,13,15\},\qquad D=A=\{0,2,4,6,8,10,12,14\},
\end{equation}
of sizes $(6,6,8,8)$, and colour every remaining line white. A cell $(r,c)$ is black-homogeneous iff $r,c\in\{5,7,9,11,13,15\}$ and $r-c,r+c$ are even; the last two conditions are automatic because $r$ and $c$ are then both odd. Hence
\[
\Bl=\{5,7,9,11,13,15\}\times\{5,7,9,11,13,15\},\qquad |\Bl|=36 .
\]
A cell is white-homogeneous iff $r,c\in\{0,1,2,3,4,6,8,10,12,14\}$ and $r-c,r+c$ are both odd, i.e.\ iff $r$ and $c$ lie in that set and have opposite parity. Writing $E=\{0,2,4,6,8,10,12,14\}$ and $F=\{1,3\}$ for its even and odd parts,
\[
\Wh=(E\times F)\cup(F\times E),\qquad |\Wh|=8\cdot2+2\cdot8=32 .
\]
By Lemma~\ref{lem:lines}, placing black queens on any $32$ of the $36$ cells of $\Bl$ and white queens on all $32$ cells of $\Wh$ gives a peaceable $32+32$ placement.
\end{proof}

Figure~\ref{fig:board16} displays the configuration, with the four omitted black-homogeneous cells taken to be the corners $(5,5),(5,15),(15,5),(15,15)$ of the $6\times6$ block.

This lower bound is not new. It is the $(4,4)$-plaid of \cite[\S3]{CDHS24}, whose peaceability is their Lemma~3.1, transported to line-set form: under the relabelling $i\mapsto i-1$ from $\{1,\dots,16\}$ to $\Zs$, their white set for $(a,b)=(4,4)$ is exactly our $\Bl$ and their black set is exactly our $\Wh$, so that the placement has $32$ queens of one colour and $36$ of the other, and $\min\{|B|,|W|\}=32$. We reproduce it here in line-set form so that Theorem~\ref{thm:main16} is self-contained and checkable by hand. Note that Theorem~1.1 of \cite{CDHS24} is asymptotic --- it asserts $t(n)\ge 0.1339\,n^2$ only for sufficiently large even $n$, and $0.1339\cdot 256>32$ --- so it is the construction, not the stated theorem, that supplies the bound at $n=16$.

\begin{figure}[t]
\centering
\setlength{\tabcolsep}{2.2pt}
\renewcommand{\arraystretch}{1.05}
{\small
\begin{tabular}{r|*{16}{c}}
$r\backslash c$ & 0 & 1 & 2 & 3 & 4 & 5 & 6 & 7 & 8 & 9 & 10 & 11 & 12 & 13 & 14 & 15 \\
\midrule
0 & $\cdot$ & $\square$ & $\cdot$ & $\square$ & $\cdot$ & $\cdot$ & $\cdot$ & $\cdot$ & $\cdot$ & $\cdot$ & $\cdot$ & $\cdot$ & $\cdot$ & $\cdot$ & $\cdot$ & $\cdot$ \\
1 & $\square$ & $\cdot$ & $\square$ & $\cdot$ & $\square$ & $\cdot$ & $\square$ & $\cdot$ & $\square$ & $\cdot$ & $\square$ & $\cdot$ & $\square$ & $\cdot$ & $\square$ & $\cdot$ \\
2 & $\cdot$ & $\square$ & $\cdot$ & $\square$ & $\cdot$ & $\cdot$ & $\cdot$ & $\cdot$ & $\cdot$ & $\cdot$ & $\cdot$ & $\cdot$ & $\cdot$ & $\cdot$ & $\cdot$ & $\cdot$ \\
3 & $\square$ & $\cdot$ & $\square$ & $\cdot$ & $\square$ & $\cdot$ & $\square$ & $\cdot$ & $\square$ & $\cdot$ & $\square$ & $\cdot$ & $\square$ & $\cdot$ & $\square$ & $\cdot$ \\
4 & $\cdot$ & $\square$ & $\cdot$ & $\square$ & $\cdot$ & $\cdot$ & $\cdot$ & $\cdot$ & $\cdot$ & $\cdot$ & $\cdot$ & $\cdot$ & $\cdot$ & $\cdot$ & $\cdot$ & $\cdot$ \\
5 & $\cdot$ & $\cdot$ & $\cdot$ & $\cdot$ & $\cdot$ & $\lozenge$ & $\cdot$ & $\blacksquare$ & $\cdot$ & $\blacksquare$ & $\cdot$ & $\blacksquare$ & $\cdot$ & $\blacksquare$ & $\cdot$ & $\lozenge$ \\
6 & $\cdot$ & $\square$ & $\cdot$ & $\square$ & $\cdot$ & $\cdot$ & $\cdot$ & $\cdot$ & $\cdot$ & $\cdot$ & $\cdot$ & $\cdot$ & $\cdot$ & $\cdot$ & $\cdot$ & $\cdot$ \\
7 & $\cdot$ & $\cdot$ & $\cdot$ & $\cdot$ & $\cdot$ & $\blacksquare$ & $\cdot$ & $\blacksquare$ & $\cdot$ & $\blacksquare$ & $\cdot$ & $\blacksquare$ & $\cdot$ & $\blacksquare$ & $\cdot$ & $\blacksquare$ \\
8 & $\cdot$ & $\square$ & $\cdot$ & $\square$ & $\cdot$ & $\cdot$ & $\cdot$ & $\cdot$ & $\cdot$ & $\cdot$ & $\cdot$ & $\cdot$ & $\cdot$ & $\cdot$ & $\cdot$ & $\cdot$ \\
9 & $\cdot$ & $\cdot$ & $\cdot$ & $\cdot$ & $\cdot$ & $\blacksquare$ & $\cdot$ & $\blacksquare$ & $\cdot$ & $\blacksquare$ & $\cdot$ & $\blacksquare$ & $\cdot$ & $\blacksquare$ & $\cdot$ & $\blacksquare$ \\
10 & $\cdot$ & $\square$ & $\cdot$ & $\square$ & $\cdot$ & $\cdot$ & $\cdot$ & $\cdot$ & $\cdot$ & $\cdot$ & $\cdot$ & $\cdot$ & $\cdot$ & $\cdot$ & $\cdot$ & $\cdot$ \\
11 & $\cdot$ & $\cdot$ & $\cdot$ & $\cdot$ & $\cdot$ & $\blacksquare$ & $\cdot$ & $\blacksquare$ & $\cdot$ & $\blacksquare$ & $\cdot$ & $\blacksquare$ & $\cdot$ & $\blacksquare$ & $\cdot$ & $\blacksquare$ \\
12 & $\cdot$ & $\square$ & $\cdot$ & $\square$ & $\cdot$ & $\cdot$ & $\cdot$ & $\cdot$ & $\cdot$ & $\cdot$ & $\cdot$ & $\cdot$ & $\cdot$ & $\cdot$ & $\cdot$ & $\cdot$ \\
13 & $\cdot$ & $\cdot$ & $\cdot$ & $\cdot$ & $\cdot$ & $\blacksquare$ & $\cdot$ & $\blacksquare$ & $\cdot$ & $\blacksquare$ & $\cdot$ & $\blacksquare$ & $\cdot$ & $\blacksquare$ & $\cdot$ & $\blacksquare$ \\
14 & $\cdot$ & $\square$ & $\cdot$ & $\square$ & $\cdot$ & $\cdot$ & $\cdot$ & $\cdot$ & $\cdot$ & $\cdot$ & $\cdot$ & $\cdot$ & $\cdot$ & $\cdot$ & $\cdot$ & $\cdot$ \\
15 & $\cdot$ & $\cdot$ & $\cdot$ & $\cdot$ & $\cdot$ & $\lozenge$ & $\cdot$ & $\blacksquare$ & $\cdot$ & $\blacksquare$ & $\cdot$ & $\blacksquare$ & $\cdot$ & $\blacksquare$ & $\cdot$ & $\lozenge$ \\
\end{tabular}}
\caption{The $32+32$ placement of Proposition~\ref{prop:lower16}: $\blacksquare$ black queens, $\square$ white queens, $\lozenge$ four of the $36$ black-homogeneous cells, left empty. Rows, columns and both diagonal families wrap around.}\label{fig:board16}
\end{figure}

\section{Parity-refined profiles for $n=16$}\label{sec:profiles16}

From now on $n=16$ and we ask whether a configuration in $\Zs$ exists with $|\Bl|\ge 33$ and $|\Wh|\ge 33$. By Lemma~\ref{lem:budget} we may assume $S\le 32$.

By Lemma~\ref{lem:incidence}(ii) the diagonal--antidiagonal incidence is no longer uniform: every cell satisfies $d\equiv a\pmod 2$, a same-parity pair of labels meets in two cells and an opposite-parity pair in none. The four sizes are therefore not enough data, and we refine them. Write
\[
r=|R|,\quad c=|C|,\quad
d_0=|D\cap 2\Z|,\quad d_1=|D\setminus 2\Z|,\quad
a_0=|A\cap 2\Z|,\quad a_1=|A\setminus 2\Z| ,
\]
and put $d=d_0+d_1$, $a=a_0+a_1$. The \emph{size profile} is the six-tuple $p=(r,c,d_0,d_1,a_0,a_1)$, with $r,c\in\{0,\dots,16\}$ and $d_0,d_1,a_0,a_1\in\{0,\dots,8\}$, and $S=r+c+d+a$.

\begin{lemma}[Pair bounds]\label{lem:pairs16}
If a configuration in $\Zs$ with profile $p$ has $|\Bl|\ge 33$ and $|\Wh|\ge 33$, then
\begin{equation}\label{eq:pair16}
rc,\ rd,\ ra,\ cd,\ ca\ \ge 33,
\qquad
2(d_0a_0+d_1a_1)\ \ge 33,
\end{equation}
and the same five products and the same parity-refined expression for the complementary sizes:
\begin{equation}\label{eq:pair16c}
\begin{gathered}
(16-r)(16-c),\ (16-r)(16-d),\ (16-r)(16-a),\\
(16-c)(16-d),\ (16-c)(16-a)\ \ge 33,\\
2\bigl((8-d_0)(8-a_0)+(8-d_1)(8-a_1)\bigr)\ \ge 33 .
\end{gathered}
\end{equation}
\end{lemma}

\begin{proof}
For the five pairs not consisting of the diagonal and antidiagonal families, Lemma~\ref{lem:incidence}(i) gives a unique common cell, so distinct cells of $\Bl$ determine distinct pairs of black lines and $|\Bl|$ is at most the product of the two sizes. For the diagonal--antidiagonal pair, Lemma~\ref{lem:incidence}(ii) gives two common cells for each same-parity pair of labels and none otherwise, so
\[
|\Bl|\ \le\ 2\,\#\{(x,y)\in D\times A: x\equiv y \pmod 2\}\ =\ 2(d_0a_0+d_1a_1).
\]
The complementary statements are the same bounds applied to $(R^c,C^c,D^c,A^c)$, whose parity-refined sizes are $(8-d_0,8-d_1,8-a_0,8-a_1)$ because each parity class of $\Zs$ has exactly $8$ labels.
\end{proof}

\begin{lemma}[Triple constraints]\label{lem:triple16}
If a configuration in $\Zs$ with profile $p$ has $|\Bl|\ge 33$ and $|\Wh|\ge 33$, then
\begin{equation}\label{eq:triple16a}
256-16(r+c+z)+rc+rz+cz\ \ge\ 66 \qquad\text{for } z\in\{d,a\},
\end{equation}
\begin{equation}\label{eq:triple16b}
256-16(x+d+a)+xd+xa+2(d_0a_0+d_1a_1)\ \ge\ 66 \qquad\text{for } x\in\{r,c\}.
\end{equation}
\end{lemma}

\begin{proof}
Apply Lemma~\ref{lem:triples} with $n=16$ and $m=33$. For the triples $\{R,C,D\}$ and $\{R,C,A\}$ all three pairs are covered by Lemma~\ref{lem:incidence}(i), so every $m(X,Y)$ is a product of sizes, giving \eqref{eq:triple16a}. For the triples $\{R,D,A\}$ and $\{C,D,A\}$ the diagonal--antidiagonal term is $2(d_0a_0+d_1a_1)$ by Lemma~\ref{lem:incidence}(ii) and the other two are products, giving \eqref{eq:triple16b}.
\end{proof}

\subsection*{Symmetry}

Three involutions act on profiles, each induced by an affine symmetry of the board:
\begin{align*}
\tau(r,c)&=(c,r): && R\leftrightarrow C,\quad D\mapsto -D,\ A\mapsto A;\\
\nu(r,c)&=(r,-c): && R\mapsto R,\ C\mapsto -C,\quad D\leftrightarrow A;\\
\sigma(r,c)&=(r+1,c): && R\mapsto R+1,\ C\mapsto C,\quad D\mapsto D+1,\ A\mapsto A+1 .
\end{align*}
Negation and translation by an even amount preserve the parity of a line label, so $\tau$ acts on profiles by $r\leftrightarrow c$ and $\nu$ by $(d_0,d_1)\leftrightarrow(a_0,a_1)$. Translation by an odd amount in one coordinate shifts both the diagonal and the antidiagonal label by an odd amount, so $\sigma$ acts by the \emph{simultaneous} exchange $d_0\leftrightarrow d_1$, $a_0\leftrightarrow a_1$. Unlike the odd case there is no halving map: $2$ is not a unit modulo $16$, the matrix $\begin{pmatrix}1&1\\1&-1\end{pmatrix}$ has determinant $-2$ and is not invertible over $\Zs$, and indeed the two pairs of families are not interchangeable, since a row and a column always meet in one cell while a diagonal and an antidiagonal meet in two or in none. The group generated by the three profile involutions has order at most $8$.

\begin{proposition}\label{prop:342}
Let $(R,C,D,A)$ be a configuration in $\Zs$ with $|\Bl|\ge 33$, $|\Wh|\ge 33$ and $S\le 32$. Then its profile is one of exactly $1{,}898$ ordered six-tuples satisfying \eqref{eq:pair16}, \eqref{eq:pair16c}, \eqref{eq:triple16a}, \eqref{eq:triple16b} and $S\le 32$, falling into exactly $342$ classes under the group generated by $\tau,\nu,\sigma$, distributed by $S$ as in Table~\ref{tab:strata16}.
\end{proposition}

\begin{table}[h]
\centering
\begin{tabular}{@{}l*{6}{r}@{}}
\toprule
$S$ & 24 & 25 & 26 & 27 & 28 & 29 \\
canonical profiles & 13 & 41 & 101 & 120 & 57 & 10 \\
\bottomrule
\end{tabular}
\medskip
\caption{The $342$ canonical parity-refined profiles for $n=16$, by budget $S$. No profile with $S\le 23$ or $S\ge 30$ survives the necessary conditions.}\label{tab:strata16}
\end{table}

\begin{proof}
Enumerate all $17^2\cdot 9^4=1{,}896{,}129$ six-tuples, filter by $S\le 32$ and by the conditions of Lemmas~\ref{lem:pairs16} and~\ref{lem:triple16}, which are necessary, and canonicalize the survivors by taking the lexicographically least element of the orbit under $\langle\tau,\nu,\sigma\rangle$. The computation is finite and was carried out independently twice, in C++ and in Python, giving $1{,}898$ ordered and $342$ canonical profiles with the distribution of Table~\ref{tab:strata16}.
\end{proof}

\begin{remark}\label{rem:no30}
Table~\ref{tab:strata16} is a structural fact worth isolating: the necessary conditions alone already exclude every profile with $S\ge 30$, so none of the top three strata $S\in\{30,31,32\}$ permitted by Lemma~\ref{lem:budget} needs to be searched at all. Contrast the odd case, where the top stratum $S=30$ contains $39$ of the $247$ canonical profiles. No single one of the conditions of Lemmas~\ref{lem:pairs16} and~\ref{lem:triple16} is responsible --- the exclusion is a joint effect and is part of the finite enumeration --- but the balanced profile at the top of the budget shows what is happening. For $n$ even and $x=y=z=n/2$, the identity \eqref{eq:triple-id} applied to $\{R,C,D\}$ gives exactly
\[
T+T'=n^2-n\cdot\tfrac{3n}{2}+3\cdot\tfrac{n^2}{4}=\tfrac{n^2}{4},
\]
whereas $|\Bl|+|\Wh|\ge 2m$ with $m=n^2/8+1$ demands $n^2/4+2$. The balanced top-stratum profile misses by exactly $2$, for every even $n$; at $n=16$ this reads $64<66$.
\end{remark}

\section{Deciding one profile exactly for $n=16$}\label{sec:completion16}

Fix a canonical profile. Since the halving map is unavailable, the fixed pair is always $(R,C)$; we may assume $r\le c$ after a transpose. The search again proceeds in three stages.

\subsection*{Stage 1: representatives of the row--column pair, and parity orientation}

Fix the pair $(R,C)$ up to the group generated by
\begin{itemize}
\item independent translations $(R,C)\mapsto(R+\alpha,\ C+\beta)$ for $\alpha,\beta\in\Zs$ ($256$ elements);
\item common unit scaling $(R,C)\mapsto(uR,\ uC)$ for odd $u$ ($8$ elements), which scales $D$ and $A$ by the same odd $u$ and hence preserves all four parity-refined sizes;
\item the transpose $\tau$, admissible exactly when $r=c$;
\item the relative sign $\nu$, i.e.\ $(R,C)\mapsto(R,-C)$, admissible \emph{only} when $(d_0,d_1)=(a_0,a_1)$, since $\nu$ exchanges the diagonal and antidiagonal families.
\end{itemize}
The group has order $256\cdot 8=2048$, doubled for each admissible extra generator.

Translations require the same care as the relative sign, and for the same reason. A translation with $\alpha-\beta$ odd sends $D\mapsto D+(\alpha-\beta)$ and $A\mapsto A+(\alpha+\beta)$ with both shifts odd, hence flips the two parity splits simultaneously: it maps the ordered profile $(r,c,d_0,d_1,a_0,a_1)$ to $(r,c,d_1,d_0,a_1,a_0)$, a different ordered profile in general. Quotienting the row--column pair by \emph{all} $256$ translations is therefore sound only if both parity orientations of every canonical profile are searched. They are: each canonical profile is expanded into the set $\{p,\sigma p\}$ of oriented profiles, which has two elements unless $p$ is $\sigma$-invariant. Across the $342$ canonical profiles this produces $677$ oriented jobs ($335$ profiles contribute two and $7$ contribute one).

Within one oriented job, the diagonal and antidiagonal families may still be exchanged, because $\nu$ does so globally, so the ordered profiles $(r,c,d_0,d_1,a_0,a_1)$ and $(r,c,a_0,a_1,d_0,d_1)$ are simultaneously feasible or infeasible; the implementation puts the family with the smaller number of admissible masks in the role of $A$, which is the family enumerated in Stage~2. The transfer argument of Lemma~\ref{lem:transfer15} applies verbatim to the group above, with its admissibility conditions ($\tau$ only when $r=c$, $\nu$ only when $(d_0,d_1)=(a_0,a_1)$): soundness requires only that the enumerated pairs meet every orbit, and the emitted set is produced by the same necklace sweep with explicit orbit marking as in the odd case. Only $14$ distinct triples $(r,c,\text{sign admissible})$ occur across the $677$ jobs, and for each of them the cardinality of the emitted set equals the orbit count computed independently by Burnside's lemma (Section~\ref{sec:verification}).

\subsection*{Stage 2: the completion as a parity-constrained subset problem}

Fix a representative $(R,C)$. For labels $x,y\in\Zs$ define
\begin{align*}
P_{x,y}&=\#\{(r,c): r-c=x,\ r+c=y,\ r\in R,\ c\in C\},\\
Q_{x,y}&=\#\{(r,c): r-c=x,\ r+c=y,\ r\notin R,\ c\notin C\} .
\end{align*}
By Lemma~\ref{lem:incidence}(ii) these are $0$ when $x\not\equiv y\pmod 2$ and lie in $\{0,1,2\}$ otherwise, the two candidate cells being $(r,c)$ and $(r+8,c+8)$ for the solution $r=\lfloor((y+x)\bmod 16)/2\rfloor$, $c=\lfloor((y-x)\bmod 16)/2\rfloor$ or its shift by $8$ in one coordinate, whichever satisfies the two original congruences.

For each antidiagonal set $A\subseteq\Zs$ with $|A\cap 2\Z|=a_0$ and $|A\setminus 2\Z|=a_1$, put
\[
p_x=\sum_{y\in A}P_{x,y},
\qquad
q_x=\sum_{y\notin A}Q_{x,y},
\qquad
Q_0=\sum_{x\in\Zs}q_x ,
\]
so that $|\Bl(R,C,D,A)|=\sum_{x\in D}p_x$ and $|\Wh(R,C,D,A)|=Q_0-\sum_{x\in D}q_x$ for every $D$. Hence a diagonal set $D$ with $|D\cap 2\Z|=d_0$ and $|D\setminus 2\Z|=d_1$ completes the configuration if and only if
\begin{equation}\label{eq:subset16}
\sum_{x\in D}p_x\ \ge\ 33
\qquad\text{and}\qquad
\sum_{x\in D}q_x\ \le\ Q_0-33 .
\end{equation}
Before \eqref{eq:subset16} is posed, the necessary conditions $\sum_x p_x\ge 33$ and $Q_0\ge 33$ are tested; they discard all but a vanishing fraction of the candidates $A$ (Table~\ref{tab:stats16}).

\subsection*{Stage 3: exact completion decision}

Problem \eqref{eq:subset16} is a $16$-item, two-constraint subset problem with a prescribed cardinality \emph{in each parity class}. It is decided in two ways, and both were run over the whole search (Section~\ref{sec:verification}).

The first is a parity-aware extremal bound followed, if the bound is inconclusive, by a depth-first search with suffix dynamic-programming bounds indexed by the two remaining cardinalities:
\begin{lemma}\label{lem:exhaustive16}
Let $\mu$ be the sum of the $d_0$ largest values $p_x$ over even $x$ plus the $d_1$ largest over odd $x$, and let $\lambda$ be the sum of the $d_0$ smallest values $q_x$ over even $x$ plus the $d_1$ smallest over odd $x$. If $\mu<33$ or $\lambda>Q_0-33$, then \eqref{eq:subset16} is infeasible. Otherwise the depth-first search decides \eqref{eq:subset16} correctly.
\end{lemma}
\begin{proof}
Any admissible $D$ selects exactly $d_0$ even and $d_1$ odd labels, so $\sum_{x\in D}p_x\le\mu$ and $\sum_{x\in D}q_x\ge\lambda$; each of the two displayed inequalities therefore contradicts one of the two conditions in \eqref{eq:subset16}. For the search, a node is abandoned only when one of the following holds: the remaining cardinality required in some parity class is negative or exceeds the number of remaining labels in that class; the accumulated $q$-sum already exceeds $Q_0-33$; the accumulated $p$-sum plus the maximal attainable remaining $p$-sum with the required parity-split cardinality is below $33$; or the accumulated $q$-sum plus the minimal attainable remaining $q$-sum with that cardinality exceeds $Q_0-33$. Each is a necessary condition for a feasible completion within the subtree, so no feasible $D$ is discarded. Success is reported only when the required cardinalities are reached and the $p$-constraint is met, or when the $p$-constraint is already met at a node whose minimal remaining $q$-sum keeps the $q$-constraint satisfiable; in both cases a feasible $D$ exists, since $p$-sums only increase as further labels are added. The search is otherwise a complete binary recursion over the $16$ labels.
\end{proof}

The second is direct enumeration: there are exactly $\binom{8}{d_0}\binom{8}{d_1}\le\binom84^2=4{,}900$ admissible masks $D$, and each is tested against \eqref{eq:subset16}. This is obviously exhaustive and involves no bound and no pruning at all.

Consequently the three stages together decide, for a given oriented profile, whether \emph{any} configuration with that profile has $|\Bl|\ge 33$ and $|\Wh|\ge 33$, and the answer transfers to the whole orbit under $\langle\tau,\nu,\sigma\rangle$.

\section{The search and its outcome for $n=16$}\label{sec:result16}

\begin{lemma}[End-to-end exhaustiveness, $n=16$]\label{lem:endtoend16}
If a configuration in $\Zs$ with $|\Bl|\ge 33$ and $|\Wh|\ge 33$ exists, then the search of Section~\ref{sec:completion16}, executed on all $677$ oriented jobs, reports a feasible completion for at least one of them.
\end{lemma}

\begin{proof}
By Lemma~\ref{lem:budget} such a configuration exists with $S\le 32$. Its parity-refined profile satisfies the conditions of Lemmas~\ref{lem:pairs16} and~\ref{lem:triple16}, so by Proposition~\ref{prop:342} some element of $\langle\tau,\nu,\sigma\rangle$ carries it to a configuration whose ordered profile is one of the $677$ oriented profiles searched; each generator is induced by a board symmetry and preserves feasibility. Within that job, a symmetry of the Stage-1 group (with the transfer argument of Lemma~\ref{lem:transfer15}) and, if applicable, the Stage-2 relabelling by $\nu$ carry the configuration to one whose row--column pair is an enumerated representative. Its antidiagonal set is among the exhaustively enumerated admissible masks, and it passes the two scalar conditions of Stage~2, which are necessary. Its diagonal set satisfies \eqref{eq:subset16}, so by Lemma~\ref{lem:exhaustive16} --- or trivially by the brute-force kernel --- the Stage-3 decision reports it feasible.
\end{proof}

\begin{proposition}\label{prop:unsat16}
For every one of the $677$ oriented profiles derived from the $342$ canonical profiles of Proposition~\ref{prop:342}, no configuration in $\Zs$ with that profile has $|\Bl|\ge 33$ and $|\Wh|\ge 33$.
\end{proposition}

\begin{proof}
The search of Section~\ref{sec:completion16} was executed twice on all $677$ jobs, by \texttt{peace16\_solver.cpp} and its brute-force variant (SHA-256 digests in Section~\ref{sec:repro}), and reported no feasible completion anywhere; by Lemma~\ref{lem:endtoend16} this establishes the claim. The run statistics are given in Table~\ref{tab:stats16}; a per-job certificate recording the oriented profile, the pair-orbit count, the number of admissible antidiagonal masks and the case counts is included among the ancillary files, for each of the two completion kernels.
\end{proof}

\begin{table}[h]
\centering
\begin{tabular}{@{}lr@{}}
\toprule
canonical profiles                          & $342$ \\
oriented profile jobs                       & $677$ \\
distinct row--column pair orbit types       & $14$ \\
Stage-2 completion subsets examined         & $13{,}163{,}028{,}768$ \\
survivors of the scalar prefilter of Stage~2 & $17{,}834$ \\
survivors of the root bounds of Stage~3 (searches entered) & $0$ \\
branch nodes in the depth-first search      & $0$ \\
jobs admitting a $33+33$ configuration      & $0$ \\
\bottomrule
\end{tabular}
\medskip
\caption{Statistics of the exhaustive run for $n=16$, one line per pruning layer as in Table~\ref{tab:stats15}. The number of Stage-2 subsets examined equals $\sum_{\text{jobs}}(\text{pair-orbit count})\cdot\binom{8}{a_0}\binom{8}{a_1}$, which is checked row by row by the audit script. The parity-aware root bound of Lemma~\ref{lem:exhaustive16} already refuted each of the $17{,}834$ scalar-prefilter survivors, so the depth-first search was never entered; the same $17{,}834$ instances were also refuted by direct enumeration of all admissible masks by the brute-force kernel.}\label{tab:stats16}
\end{table}

The collapse from $1.3\cdot10^{10}$ candidates to $17{,}834$ is caused entirely by the two necessary conditions $\sum_x p_x\ge 33$ and $Q_0\ge 33$ of Stage~2; both are conditions on the antidiagonal set alone, and on the even torus they are severe.

\begin{proof}[Proof of Theorem~\ref{thm:main16}]
The lower bound is Proposition~\ref{prop:lower16}. For the upper bound, suppose $33$ black and $33$ white queens are placed peaceably. By Lemma~\ref{lem:lines} there is a configuration with $|\Bl|\ge 33$ and $|\Wh|\ge 33$, by Lemma~\ref{lem:budget} one may take $S\le 32$, and by Proposition~\ref{prop:342} its parity-refined profile lies in the $342$-element canonical list, hence its ordered profile is one of the $677$ oriented profiles searched. This contradicts Proposition~\ref{prop:unsat16}. Hence $t(16)\le 32$, and $t(16)=32$.
\end{proof}

\section{Verification}\label{sec:verification}

An exhaustive computation is only as good as the evidence that it is exhaustive and correctly implemented. We list the checks performed, case by case. The evidence is not equally strong on the two sides, and the difference is stated plainly in Section~\ref{sec:ver16}.

\subsection{What both cases have in common}\label{sec:vercommon}

\emph{Startup gates.} Each solver begins with gates that abort the run on failure rather than proceeding: a randomized comparison of its exact completion decision against direct enumeration of all admissible masks, a check of known orbit counts, and a regeneration of the profile list from the necessary conditions.

\emph{Memory safety.} AddressSanitizer and UndefinedBehaviorSanitizer builds of both proof solvers pass all startup self-tests; beyond that, an ASan+UBSan build of each proof solver re-ran its \emph{complete} sweep --- all $247$ profiles for $n=15$ ($232$~s single-threaded) and all $677$ jobs for $n=16$ ($244$~s) --- with zero sanitizer diagnostics and certificates identical to the recorded runs in every non-timing column.

\emph{Independent audit scripts.} A Python audit script for each case, sharing no code with the solvers, re-derives the profile list, recomputes every fixed-pair orbit count from Burnside's lemma --- counting, for each group element, the subsets of each cardinality invariant under the induced affine permutation of $\Zn$, via the cycle structure --- checks the per-row and total case arithmetic of the certificates, and re-counts the lower-bound witness. Both print \texttt{AUDIT\_OK}.

\emph{What is not claimed.} The certificate files are audit logs, not standalone proof objects in the sense of DRAT or LRAT: they record what was searched, not a machine-checkable refutation. The refutation is the short program together with the reductions of Sections~\ref{sec:reduction}--\ref{sec:completion16}; re-running the program is what re-establishes it, and agreement between implementations is a strong error filter, not a formal guarantee.

\subsection{The case $n=15$}\label{sec:ver15}

The two enumerators mentioned below --- the \emph{proof solver} and the \emph{independent enumerator} --- were written separately, from the mathematical description of Section~\ref{sec:completion15}, in different C++ dialects, by different systems (see Section~\ref{sec:ai}), and share no code.

\emph{Two independent enumerators.} The proof solver (\texttt{peace15\_solver.cpp}) and the independent enumerator (\texttt{profile\_enum.cpp}) agree that all $247$ profiles are infeasible. They differ in three substantive respects. First, in the symmetry quotient: the two select different fixed pairs on $19$ of the $247$ profiles, and consequently enumerate different case sets --- $6{,}074{,}753{,}568$ antidiagonal subsets for the solver against $6{,}241{,}793{,}402$ completions for the enumerator. Second, in the decision procedure for \eqref{eq:subset15}: the solver uses the depth-first search with suffix dynamic-programming bounds of Stage~3, whereas the enumerator uses an exact $7+8$ meet-in-the-middle table, retaining for each cardinality and each $p$-sum threshold the minimum achievable $q$-sum over right-half subsets --- an algorithm with no search and no pruning at all. Third, in scheduling: the solver is OpenMP-parallel, the enumerator single-threaded and resumable, writing one JSON record per profile.

\emph{Agreement on the profile list.} The $247$-profile worklist was generated independently by us in Python, from Lemmas~\ref{lem:pairs15} and~\ref{lem:triples}; the proof solver regenerates the set from scratch in C++ at startup and aborts unless it coincides with the supplied file, which it does byte for byte; the audit script regenerates it a third time in Python. At $S=30$ the count was also reproduced by an independent hand-checkable enumeration as $370$ ordered profiles falling into $39$ canonical classes.

\emph{Repeated runs on different toolchains.} The proof solver was run three times: twice as recorded by its author (five threads, $60.4$~s; three threads, $74.0$~s) and once by us from a fresh build with a different compiler and no OpenMP (single-threaded, $101.7$~s). All $247$ rows of the certificate agree in every column except the per-profile timing.

\emph{Brute-force gates on the completion kernel.} At startup the proof solver compares its Stage~3 decision against direct enumeration of all $\binom{15}{k}$ subsets on $20{,}000$ deterministic pseudorandom instances with randomized weights and cardinalities, and aborts on any disagreement; the enumerator runs the same style of gate on $100{,}000$ instances. The solver additionally checks three known pair-orbit counts before starting: $11{,}793$ for $(r,c)=(7,7)$ under the order-$3600$ group (transpose, no relative sign), $6{,}892$ for $(7,7)$ under the order-$7200$ group (transpose and relative sign), and $13{,}654$ for $(7,8)$ with the relative sign; the first two had been derived independently earlier in the project, by Burnside's lemma.

\emph{Burnside audit.} The audit script (\texttt{peace15\_audit.py}) recomputes all $43$ fixed-pair orbit counts from Burnside's lemma and compares them with the counts recorded in the certificate. It also verifies that each certificate row satisfies $\text{cases}=\text{orbits}\times\binom{15}{a}$, verifies the column total $6{,}074{,}753{,}568$, and re-counts the witness of Proposition~\ref{prop:lower15} as $20$ black-homogeneous and $22$ white-homogeneous cells.

\emph{The witness.} The $20+20$ placement of Proposition~\ref{prop:lower15} was checked by four independently written scripts, including a from-scratch check that verifies all $400$ black--white pairs against all four attack relations directly, rather than through the line-set formulation. That check is included among the ancillary files as \texttt{check\_witness15.py}; the other three (the audit script's recount, and two checks in the project repository) are described here for completeness.

\emph{Source audit.} The proof solver was read line by line against the mathematics of Section~\ref{sec:completion15}, with particular attention to the three places where an unsound shortcut would be invisible in the output: the admissibility conditions for the transpose and relative-sign quotients, the necessary-condition character of every pruning rule, and the reconstruction path that would turn a reported feasible instance into an explicit quadruple verified cell by cell.

\emph{Independent partial evidence.} A SAT-based cube-and-conquer campaign on a $3{,}244$-variable encoding of the $21+21$ question, later superseded and never finished, produced roughly $28{,}900$ DRAT-certified cube refutations checked by \texttt{drat-trim} in default backward mode; none of the processed cubes was satisfiable, though the cube partition was never completed. It proves nothing on its own but is consistent with Theorem~\ref{thm:main15}. A per-profile certified-CNF pilot refuted $(5,5,5,5)$ with a checked $9.2$~MB DRAT proof but timed out on $(6,8,7,7)$ under two solvers; this is why an exhaustive enumeration, and not a proof-logging pipeline, carries the upper bound.

\subsection{The case $n=16$}\label{sec:ver16}

\emph{The gap in the evidence, stated first.} For $n=16$ there is as yet \emph{no fully independent second enumerator}. Two programs were written and run to completion, \texttt{peace16\_solver.cpp} and \texttt{peace16\_solver\_bruteforce.cpp}, but they are the same program apart from the completion kernel: profile generation, the symmetry quotient, the pair-representative construction, the parity orientation logic and the antidiagonal loop are shared verbatim, and the two files differ in ten lines. A systematic error in the shared part --- for instance an unsound quotient that discarded a feasible configuration before the kernel ever saw it --- would not be caught by their agreement. This is a real weakness relative to the $n=15$ case, and closing it by writing a genuinely independent enumerator from Sections~\ref{sec:profiles16}--\ref{sec:completion16} is the single most valuable check that remains open. What follows should be read against that reservation.

\emph{Two completion kernels over the whole search.} The two programs agree on every non-timing column of all $677$ certificate rows: the same $13{,}163{,}028{,}768$ antidiagonal subsets examined, the same $17{,}834$ survivors of the prefilter, and \texttt{UNSAT} for every job. Where they differ, they differ substantively. The optimized kernel refuted each of the $17{,}834$ surviving instances by the parity-aware extremal bound of Lemma~\ref{lem:exhaustive16} alone, without entering its depth-first search --- the certificates record $0$ branch nodes. The brute-force kernel refuted the same $17{,}834$ instances by testing every one of the $\binom{8}{d_0}\binom{8}{d_1}\le 4{,}900$ admissible masks. So the instances that actually decide the theorem were settled twice, once by a counting bound and once by complete enumeration, with no shared code on that path. Each binary independently passes a startup gate comparing the two kernels on $20{,}000$ deterministic pseudorandom instances with randomized weights and parity cardinalities.

\emph{Independent audit.} The audit script \texttt{peace16\_audit.py} shares no code with the solvers and re-derives, in Python: the $1{,}898$ ordered and $342$ canonical profiles and their strata from the necessary conditions of Lemmas~\ref{lem:pairs16} and~\ref{lem:triple16}; all $14$ row--column pair-orbit counts by Burnside's lemma; the two-lift reconstruction of Lemma~\ref{lem:incidence}(ii), by comparing the closed-form formula used in the solver against brute-force enumeration of the common cells for all $256$ label pairs $(d,a)$; the lower-bound witness \eqref{eq:witness16}, recounted as $36$ black-homogeneous and $32$ white-homogeneous cells; the row arithmetic and the column totals $13{,}163{,}028{,}768$ and $17{,}834$ of both certificates, and their column-by-column agreement. It was re-run locally during preparation of this note and prints \texttt{AUDIT\_OK}.

\emph{Startup gates in the solvers.} Beyond the $20{,}000$-instance kernel comparison, each solver checks that the number of $8$-element necklaces in $\Zs$ is $810$ and that two known pair-orbit counts, $73{,}663$ for $(r,c)=(7,8)$ and $42{,}062$ for $(8,8)$ without the relative sign, come out right; it regenerates the profile list and aborts unless it finds $1{,}898$ ordered and $342$ canonical profiles. Had any job been satisfiable, the solver would have reconstructed an explicit quadruple $(R,C,D,A)$ and verified it by direct enumeration of the $256$ cells before reporting; that path is exercised by the self-tests and was never taken in the run.

\emph{Local rebuild and rerun.} During preparation of this note both
programs were additionally rebuilt with a different toolchain (Apple
\texttt{clang} on \texttt{arm64}, no OpenMP, single-threaded) and rerun to
completion; all $677$ rows of both regenerated certificates agree with the
recorded five-thread runs in every non-timing column. The startup gates of
the preceding paragraph were exercised by these reruns.

\emph{The witness.} The $32+32$ placement of Proposition~\ref{prop:lower16} was verified independently of the line-set formulation by checking all $32\cdot 32=1024$ black--white pairs against all four attack relations directly (included among the ancillary files as \texttt{check\_witness16.py}), and separately by reconstructing the $(4,4)$-plaid from the definition in \cite[\S3]{CDHS24} and confirming that it coincides cell for cell with \eqref{eq:witness16} after the relabelling $i\mapsto i-1$.

\emph{Source audit.} The proof solver was read against the mathematics of Sections~\ref{sec:profiles16}--\ref{sec:completion16}, with particular attention to the two soundness points peculiar to the even case: that quotienting the row--column pair by all $256$ translations is compensated by searching both parity orientations of every canonical profile, and that the relative-sign quotient is applied only when $(d_0,d_1)=(a_0,a_1)$.

\section{Reproduction and ancillary files}\label{sec:repro}

The ancillary files accompanying this note are:
\begin{center}
\begin{tabular}{@{}ll@{}}
\toprule
\multicolumn{2}{@{}l}{\emph{$n=15$}}\\
\texttt{peace15\_solver.cpp} & the proof solver (C++17, OpenMP)\\
\texttt{profile\_enum.cpp} & the independent enumerator (C++20)\\
\texttt{peace15\_audit.py} & profile, Burnside and witness audit (Python 3)\\
\texttt{WORKLIST-2026-07-25.tsv} & the $247$ canonical profiles\\
\texttt{peace15\_certificate.tsv} & per-profile certificate, five-thread run\\
\texttt{peace15\_certificate\_run2.tsv} & per-profile certificate, three-thread run\\
\addlinespace
\multicolumn{2}{@{}l}{\emph{$n=16$}}\\
\texttt{peace16\_solver.cpp} & the proof solver (C++20, OpenMP)\\
\texttt{peace16\_solver\_bruteforce.cpp} & the same, with the brute-force kernel\\
\texttt{peace16\_audit.py} & profile, Burnside, lift-map and witness audit\\
\texttt{peace16\_certificate.tsv} & per-job certificate, optimized kernel\\
\texttt{peace16\_certificate\_bruteforce.tsv} & per-job certificate, brute-force kernel\\
\addlinespace
\texttt{check\_witness15.py} & pairwise check of the $20+20$ witness\\
\texttt{check\_witness16.py} & pairwise check of the $32+32$ witness\\
\texttt{README.txt} & instructions, certificate schema, expected output\\
\texttt{SHA256SUMS} & digests of every other ancillary file\\
\bottomrule
\end{tabular}
\end{center}

To rebuild the upper bound for $n=15$ (portable verification build):
\begin{verbatim}
g++ -O3 -fopenmp -std=c++17 peace15_solver.cpp -o peace15_solver
OMP_NUM_THREADS=5 ./peace15_solver \
    --worklist WORKLIST-2026-07-25.tsv \
    --output peace15_certificate.tsv
python3 peace15_audit.py
\end{verbatim}
The solver prints \texttt{SELFTEST\_OK} after its startup gates and \texttt{ALL\_UNSAT} on completion; the audit script prints \texttt{AUDIT\_OK}. The recorded five-thread run took about $60$ seconds. The independent enumerator is built and run by
\begin{verbatim}
clang++ -O3 -std=c++20 profile_enum.cpp -o profile_enum
./profile_enum --self-test --random 100000
./profile_enum --batch --worklist WORKLIST-2026-07-25.tsv \
    --results results.jsonl --log logs
\end{verbatim}
and takes about $200$ seconds single-threaded.

To rebuild the upper bound for $n=16$ (portable verification build):
\begin{verbatim}
g++ -O3 -fopenmp -std=c++20 peace16_solver.cpp -o peace16_solver
OMP_NUM_THREADS=5 ./peace16_solver --output peace16_certificate.tsv

g++ -O3 -fopenmp -std=c++20 \
    peace16_solver_bruteforce.cpp -o peace16_solver_bruteforce
OMP_NUM_THREADS=5 ./peace16_solver_bruteforce \
    --output peace16_certificate_bruteforce.tsv

python3 peace16_audit.py
\end{verbatim}
Each solver prints \texttt{SELFTEST\_OK} after its startup gates and \texttt{ALL\_UNSAT} on completion; the audit script prints \texttt{AUDIT\_OK} and compares the two certificates column by column. The recorded five-thread runs took about $78$ seconds each. Runtime is not part of either proof. Note that both audit scripts read the certificates from their own directory, so run them after the solvers have written those files, or against the shipped copies.

Three remarks on the build and run environment. First, \texttt{-march=native} and \texttt{-DNDEBUG} may be added for speed and were used in the recorded runs; they are omitted from the portable commands above because neither affects any verification step --- in particular, no proof-relevant check in any of the programs is implemented via \texttt{assert} (the only assertion is an internal size-ordering invariant), all gates being explicit runtime tests that abort the program and are active in every build. Second, the randomized self-tests are deterministic: they use \texttt{std::mt19937\_64} with the fixed seeds \texttt{0x4d595df4d0f33173} ($n=15$, $20{,}000$ instances) and \texttt{0x16e0ddc0ffee} ($n=16$, $20{,}000$ instances per solver), so every rerun tests the same instances. Third, the recorded five-thread runs were produced with GCC/OpenMP in an x86-64 (AMD EPYC) container, and the independent local reruns --- serial, and under AddressSanitizer/UndefinedBehaviorSanitizer --- with Apple \texttt{clang}~21 on an \texttt{arm64} (Apple Silicon) machine running macOS; agreement of all non-timing certificate columns across these environments is part of the evidence of Section~\ref{sec:verification}.

SHA-256 digests of the files that carry the two upper bounds, as recorded:
\begin{center}\small
\begin{tabular}{@{}l@{}}
\toprule
\texttt{peace15\_solver.cpp}\\
\quad\texttt{1ae5b340d6feb4591dae4f7b57a283a373ce60fd2bc9767fc043f25d402bebe9}\\
\addlinespace
\texttt{peace15\_certificate.tsv}\\
\quad\texttt{f9d2c858cc6e3a144a10da509a744114080de801df6d216b8cd2783cb2f410b0}\\
\addlinespace
\texttt{peace16\_solver.cpp}\\
\quad\texttt{ff8de5710446f1b8721e372a3f9af70bcd05b5798cb7924bce58d7e9c784df7c}\\
\addlinespace
\texttt{peace16\_solver\_bruteforce.cpp}\\
\quad\texttt{4153d2181642b3de83f3f0f366d85b8438350aa3b620580cc3509fafe7c65c3f}\\
\addlinespace
\texttt{peace16\_certificate.tsv}\\
\quad\texttt{850fc499127193539c06216aa57e91fd6373b3ea2773f4c642654e95def0c086}\\
\addlinespace
\texttt{peace16\_certificate\_bruteforce.tsv}\\
\quad\texttt{8f606eac6e0b3bcc31f07b21bcfabf1fddd5d6abd350c9b557765f01dc2948d2}\\
\bottomrule
\end{tabular}
\end{center}

\section{Use of AI tools}\label{sec:ai}

This note was produced with substantial use of AI systems, which we disclose in the interest of transparency. The reduction programme --- the line-colouring equivalence, the budget inequality, the pair and triple bounds, the symmetry analysis and the $247$-profile worklist --- was developed inside an AI-assisted research repository in which Anthropic Claude agents ran the computations and maintained the standing documents, and in which a large-language-model collaborator acting as a mathematician was consulted in writing; the exact completion algorithm of Section~\ref{sec:completion15} originates in that consultation. One error in that reply (a claim that the budget could be normalized to $S=2n$) was found and corrected here; see Remark~\ref{rem:notequal}. The final proof documents for both cases, all three proof solvers and both audit scripts listed in Section~\ref{sec:repro} were produced by OpenAI GPT-5.6~Pro on July~25, 2026.

The verification described in Section~\ref{sec:verification} was carried out by Anthropic Claude agents: for $n=15$, the second exact enumerator \texttt{profile\_enum.cpp} and its brute-force gates, the line-by-line audit of the proof solver, the local rebuild and rerun on a different toolchain, the sanitizer builds, the re-derivation and numerical checking of every identity in Sections~\ref{sec:reduction}--\ref{sec:profiles15}, and the independent checks of the witness; for $n=16$, the re-derivation of the identities of Sections~\ref{sec:profiles16}--\ref{sec:completion16}, the source audit of the solver, the independent recount of the witness and of the $(4,4)$-plaid, and the local re-run of the audit script. The soundness correction to the fixed-pair quotient described in Section~\ref{sec:completion15} was found during that verification.

Agreement between AI-written implementations is a useful and, for $n=15$, deliberately engineered error filter, but it is not independent verification in the scholarly sense, and it does not replace the proofs given in the text. For $n=16$ even that filter is weaker, as Section~\ref{sec:ver16} states. Readers who wish to check the theorems should re-run the ancillary programs or, better, write a third implementation from Sections~\ref{sec:reduction}--\ref{sec:completion16}. The author reviewed the manuscript and takes full responsibility for its content.

\section*{Acknowledgements}

The problem was brought within reach by the work of Clinch, Drescher, Huynh and Saffidine, whose odd-torus bound bracketed $t(15)$ from above, whose even-torus bound bracketed $t(16)$, and whose constructions supply both matching lower bounds, and by the OEIS contributors who have kept A279405 current --- most recently Jaideep Padhi, whose values for $n=13$ and $n=14$ in June 2026 marked $n=15$ as the live frontier.

\begin{thebibliography}{99}

\bibitem{Ainley77}
S.~Ainley,
\emph{Mathematical Puzzles},
G.~Bell \& Sons, London, 1977.

\bibitem{CDHS24}
K.~Clinch, M.~Drescher, T.~Huynh, and A.~Saffidine,
\emph{Constructions, bounds, and algorithms for peaceable queens},
arXiv:2406.06974v3 (2024).

\bibitem{OEIS250000}
N.~J.~A. Sloane and others,
\emph{The On-Line Encyclopedia of Integer Sequences}, sequence A250000:
\emph{Peaceable coexisting armies of queens},
\url{https://oeis.org/A250000}, accessed July 25, 2026.

\bibitem{OEIS279405}
N.~J.~A. Sloane and others,
\emph{The On-Line Encyclopedia of Integer Sequences}, sequence A279405:
\emph{Peaceable coexisting armies of queens on a toroidal board},
\url{https://oeis.org/A279405}, accessed July 25, 2026.

\end{thebibliography}

\end{document}
